\begin{figure*}[!htb]
    \centering
    \includegraphics[width=\linewidth]{./Figure/PhaseWidth.png}
    \caption{
        (a) Phase width variation for different sliding side-chain fractions, $\phi_{\mathrm a}=0.25$, $0.50$, and $0.75$.
        The upper line plot compares the normalized widths of the L, G, C, BCC, and FCC regions at $N\chi=40.0$, defined as $\Delta\bar{\phi}_{\mathrm A}(\phi_{\mathrm a})/\Delta\bar{\phi}_{\mathrm A}(0.25)$.
        The lower stacked bars show the phase widths along the $\bar{\phi}_{\mathrm A}$ direction for different $\phi_{\mathrm a}$ at $N\chi=25.0$, $35.0$, and $45.0$.
        (b--d) BCC phase region.
        (b) Width of the BCC region, $\Delta\bar{\phi}_{\mathrm A}$, as a function of $N\chi$ for $\phi_{\mathrm a}=0.25$, $0.50$, and $0.75$.
        C and S denote the A-core and A-shell composition ranges, respectively; circles and triangles denote the corresponding ranges, and dashed lines are spline fits.
        (c) Isodensity layer preferred topology $f_{\Omega}^{\ast}$ as a function of $\phi_{\mathrm A}^{\mathrm{iso}}$ at different $N\chi$ for fixed $\phi_{\mathrm a}=0.50$.
        (d) $f_{\Omega}^{\ast}$ as a function of $\phi_{\mathrm a}$ at fixed $\phi_{\mathrm A}^{\mathrm{iso}}=0.75$ for different $N\chi$.
        (e--g) Same as (b--d), but for the FCC phase.
    }
    \label{fig:PhaseWidth}
\end{figure*}


% 综上所述，滑动自由度 $f$ 提供了额外的构象调节机制，使 L-S TT 能够在不同空间位置自发调整拓扑，以调整自发曲率和适应有序相中的局部密度要求，缓解 packing frustration。
% 从自由能角度看，$f$ 的作用主要体现在两个方面。  
% 其一，$f$ 的引入扩展了体系可访问的拓扑态空间，从而带来由滑动自由度产生的额外熵贡献。  
% 其二，$f$ 的空间重分布会改变链在界面附近的伸展方式，使链具有了对局域密度匹配以及对曲率的适应能力。
% 这些由 $f$ 引入所带来的影响，最终都体现在自由能中熵项与焓项的竞争之中。
% 随着 $N\chi$ 增大，界面相关的焓贡献逐渐增强，因此滑动熵与局域焓优化这两种效应的相对权重也会发生变化，进而使 L-S TT 在低 $N\chi$ 下表现出拓扑择优，对应低滑动熵，而在高 $N\chi$ 下转而表现出拓扑均化，对应高滑动熵。
% 具体地，在低 $N\chi$ 区域，体系对构象调节的需求相对较弱，此时仅出现层状相及双联通相等界面形状较为简单、曲率变化较缓的结构，局域环境对拓扑的要求也相对一致，因而体系倾向于围绕某一择优的拓扑 $f_\Omega^\ast$ 形成单峰分布。
% 在这一条件下，SCFT 得到的最优拓扑与 RPA 单模近似预测的 $f_\Omega^\ast$ 近似一致，拓扑涨落所对应的额外构象熵主要体现在单峰分布宽度的变化上。
% 当 $N \chi$ 增大且体系形成的有序结构具有复杂界面（如 BCC 与 FCC 相）时，调节 $f$ 能显著降低界面能与链段伸展挫折所带来的能量代价。
% 此时，体系不再偏向单一的择优拓扑，而是呈现出更宽的 $f$ 分布范围；不同拓扑构型在魏格纳–赛茨元胞的不同空间位置上发生功能分化，并协同分担局域填充任务。
% 另一方面，即便在高 $N \chi$ 条件下，若目标相结构界面形貌较为规整（例如图~\ref{fig:Gyroid} 中的 Gyroid 相），调节 $f$ 带来的界面能优化有限，系统仍可能维持受构象熵主导的$f$分布。
In summary, the sliding degree of freedom $f$ provides an additional mechanism for conformational regulation, allowing L--S TT copolymers to redistribute the contributions of different topologies to the local A-segment density in response to local interfacial curvature and filling environments.
This local topological adaptation helps accommodate structural requirements in ordered phases and may thereby relieve packing frustration.


Allowing $f$ to vary expands the accessible molecular configurations and enables the topology weights to redistribute self-consistently in space.
The resulting topology redistribution modifies the chain conformational response near interfaces, coupling the conformational and sliding distribution entropy responses to local density matching and curvature response through the self-consistent fields.



As $N\chi$ increases, stronger A--B repulsion enhances the interfacial and packing constraints, changing the balance between conformational regulation and local geometric adaptation.

The observed evolution from pronounced topology selection at low $N\chi$ toward broader topology distributions at high $N\chi$ reflects this changing balance.

Specifically, in the low-$N\chi$ regime, the need for conformational adjustment is relatively weak.
Consistently, the stable ordered phases in this regime are mainly those with relatively simple interfacial geometry and low interfacial curvature, such as L and G phases.
Because the local environments impose similar topological requirements, the system tends to form a single-peaked distribution centered around a preferred topology $f_\Lambda^\ast$.

Under these conditions, the preferred topology obtained from SCFT is approximately consistent with $f_{\mathrm{RPA\text{-}MUT}}^\ast$ predicted by RPA--MUT.
The width of the single-peaked distribution characterizes the range of contributing topologies, while the associated conformational response is contained in the chain statistics.

When $N\chi$ increases and the ordered structures develop more complex interfaces, such as in the BCC and FCC phases, adjusting $f$ provides an additional route for accommodating interfacial geometry and heterogeneous local filling requirements.
The system then no longer favors a single preferred topology, but instead develops a broader distribution of $f$.
Different topological configurations become functionally differentiated at different positions within the Wigner--Seitz cell and cooperatively accommodate local filling requirements.

Even at high $N\chi$, an ordered structure with relatively regular interfacial geometry, such as the Gyroid phase in Fig.~\ref{fig:Gyroid}, may retain a pronounced topology preference.
This behavior suggests that the local geometric requirements favor similar topological configurations across different spatial regions.



% 通过比较不同 $N\chi$ 条件下相区宽度的变化，可以揭示体系中熵效应与焓效应之间的竞争关系。
% 图~\ref{fig:PhaseWidth}a 展示了在不同 $\phi_\mathrm{a}$ 下，L-S TT 体系随 $N\chi$ 增加所形成的有序相区宽度的变化趋势。
% 柱状图给出了不同相在总有序区宽度中的占比变化；上方折线图则标准化比较了各 $\phi_\mathrm{a}$ 条件下有序区宽度 $\Delta \bar{\phi}_{\mathrm{A}}(\phi_\mathrm{a},X)$ 的相对变化。
% 其中，$X \in \{\mathrm{L},\mathrm{G},\mathrm{C},\mathrm{BCC},\mathrm{FCC}\}$ 表示不同的有序相结构。
% 整体来看，随着 $N\chi$ 的增加，总有序相区的宽度呈增长趋势，符合嵌段共聚物在组分不相容性增强下，自发形成有序结构以降低焓代价的普遍特征。
% 然而，该增长趋势在不同 $\phi_\mathrm{a}$ 条件下的表现并不一致，这主要源于两方面机制的耦合：构象熵对 $\phi_\mathrm{a}$ 依赖的非线性特征与不同相所对应界面结构复杂度的差异。
% 首先，图~\ref{fig:ODT}c 已表明，L-S TT 对构象熵的调节能力并不随 $\phi_\mathrm{a}$ 单调增强，而是呈现出先增后降的非单调变化\fc{这里把长侧链和短侧链的论述删掉了，前文（第四节开始第一段，以及对滑动熵部分的讨论）已经讨论过这一结果，这里应该不用再重复叙述一遍。}。
% 相应地，在中等滑动比例 $\phi_\mathrm{a}\approx 0.50$ 附近，体系对滑动自由度的利用最为充分，因而其构象熵调节能力也最强。
% 其次，各有序相的界面结构复杂程度不同，对滑动调节能力的依赖也各异，从而导致不同相在 $\phi_\mathrm{a}$ 增加过程中的响应趋势有所不同。
% L 相具有规则平直界面，但要求 A、B 组分的体积比接近，而 L-S TT 一般不满足该条件，构象调节对其稳定性的增益有限，故其占比表现出先减后增的非单调行为，如图~\ref{fig:PhaseWidth}a 所示。
% G 相包含互穿通道，对链段的柔性调节更敏感，因此其相区占比随 $\phi_\mathrm{a}$ 增大而显著下降。
% 相比之下，C 相以及 BCC 与 FCC 等球状相结构复杂、界面高度弯曲且填充紧密，对滑动自由度调控密度场的能力依赖更强，因而其相区宽度在 $\phi_\mathrm{a} \approx 0.50$ 附近达到最大值。
% 对 $\phi_\mathrm{a}$ 依赖的差异性在相区宽度增幅比中同样存在，比较图~\ref{fig:PhaseWidth}a 中 $\phi_{\mathrm{a}}=0.25$ 和 $\phi_\mathrm{a}=0.50$ 的结果可得：
% 这种对 $\phi_\mathrm{a}$ 的在相区宽度增幅比中同样存在。
% 由图~\ref{fig:PhaseWidth}a 可见，当以 $\phi_\mathrm{a}=0.25$ 为基准进行比较时，$\phi_\mathrm{a}=0.50$ 的结果表明：
% \begin{equation}
%     \begin{split}
%         \frac{\Delta \bar{\phi}_{\mathrm{A}}(\phi_{\mathrm{a}}=0.50,\mathrm{C})}
%         {\Delta \bar{\phi}_{\mathrm{A}}(\phi_{\mathrm{a}}=0.25,\mathrm{C})} <&
%         \frac{\Delta \bar{\phi}_{\mathrm{A}}(\phi_{\mathrm{a}}=0.50,\mathrm{BCC})}
%         {\Delta \bar{\phi}_{\mathrm{A}}(\phi_{\mathrm{a}}=0.25,\mathrm{BCC})} \\
%         <&\frac{\Delta \bar{\phi}_{\mathrm{A}}(\phi_{\mathrm{a}}=0.50,\mathrm{FCC})}
%         {\Delta \bar{\phi}_{\mathrm{A}}(\phi_{\mathrm{a}}=0.25,\mathrm{FCC})} 
%     \end{split}
% \end{equation}
% 该不等式表明，滑动自由度对界面复杂度较高相的稳定化作用更强，且其对 $\phi_{\mathrm{a}}$ 变化的响应也更为显著；特别是在 $\phi_\mathrm{a}\approx 0.50$ 附近，这种作用在 FCC 相中最为突出，表明在 $N\chi$ 较高时，L-S TT 有利于球状相的形成。
By comparing the variation of phase widths at different $N\chi$, we can clarify the competition between entropic and enthalpic effects in the L--S TT system.
Fig.~\ref{fig:PhaseWidth}a shows how the widths of the ordered phase regions change with increasing $N\chi$ for different $\phi_{\mathrm a}$.
The stacked bar plot gives the relative contribution of each phase to the total ordered-region width, whereas the upper line plot compares the normalized width $\Delta\bar{\phi}_{\mathrm A}(\phi_{\mathrm a},X)$ for different $\phi_{\mathrm a}$.
Here, $X\in\{\mathrm{L},\mathrm{G},\mathrm{C},\mathrm{BCC},\mathrm{FCC}\}$ denotes the ordered phase.
Overall, the total width of the ordered region increases with $N\chi$, consistent with the general tendency of block copolymers to form ordered structures when the incompatibility between blocks becomes stronger.
However, this increase depends strongly on $\phi_{\mathrm a}$.

We interpret this dependence through the coupling between the conformational response to molecular architecture and the interfacial requirements of the ordered phases.
The effective conformational-entropy scale in Fig.~\ref{fig:ODT}e varies nonmonotonically with $\phi_{\mathrm a}$, providing a reference for this interpretation rather than a direct entropy decomposition for each ordered phase.

The response to $\phi_{\mathrm a}$ also depends on the interfacial geometry of each ordered phase.
The L phase has regular flat interfaces, but it requires nearly balanced A and B volume fractions, a condition not generally satisfied by L--S TT copolymers; therefore, conformational regulation provides only a limited stabilization effect, leading to a nonmonotonic change in its phase fraction, as shown in Fig.~\ref{fig:PhaseWidth}a.
The G phase contains interpenetrating network channels and is more sensitive to flexible chain regulation, so its phase fraction decreases markedly as $\phi_{\mathrm a}$ increases.
By contrast, the C phase and the spherical phases, including BCC and FCC, have more complex and highly curved interfaces and more constrained packing environments.
They therefore rely more strongly on the ability of the sliding degree of freedom to regulate the density field, and their phase widths reach maxima near $\phi_{\mathrm a}\approx0.50$.
This phase-dependent response is also reflected in the relative increase of phase widths.
Using $\phi_{\mathrm a}=0.25$ as the reference, the results at $\phi_{\mathrm a}=0.50$ in Fig.~\ref{fig:PhaseWidth}a show that
\begin{equation}
    \begin{split}
        \frac{\Delta \bar{\phi}_{\mathrm{A}}(\phi_{\mathrm{a}}=0.50,\mathrm{C})}
        {\Delta \bar{\phi}_{\mathrm{A}}(\phi_{\mathrm{a}}=0.25,\mathrm{C})} <
        \frac{\Delta \bar{\phi}_{\mathrm{A}}(\phi_{\mathrm{a}}=0.50,\mathrm{BCC})}
        {\Delta \bar{\phi}_{\mathrm{A}}(\phi_{\mathrm{a}}=0.25,\mathrm{BCC})} 
        <
        \frac{\Delta \bar{\phi}_{\mathrm{A}}(\phi_{\mathrm{a}}=0.50,\mathrm{FCC})}
        {\Delta \bar{\phi}_{\mathrm{A}}(\phi_{\mathrm{a}}=0.25,\mathrm{FCC})}.
    \end{split}
\end{equation}
This inequality indicates that the stabilizing effect of the sliding degree of freedom becomes stronger for phases with more complex interfacial geometry.
The effect is strongest for the FCC phase, suggesting that at relatively high $N\chi$, L--S TT copolymers favor the formation of spherical phases.


% 为了阐明不同拓扑调节机制在复杂有序结构形成中的作用，我们考察了图~\ref{fig:PhaseWidth}b,e 所示的 BCC 与 FCC 相区宽度 $\Delta \bar{\phi}_{\mathrm{A}}$ 随 $N\chi$ 和 $\phi_\mathrm{a}$ 的演化行为。
% 结果表明，在所有 $\phi_\mathrm{a}$ 下，BCC 相区宽度均随 $N\chi$ 单调增加，这说明随着焓的增加，BCC 相始终能够稳定受益于滑动自由度$f$的引入所带来的构象调节。
% 结合不同 $\phi_\mathrm{a}$ 条件下各等值面 $\phi_{\mathrm A}^{\mathrm{iso}}$ 上的 $f_\Omega^\ast$ 的分布特征可解释该现象。
% 如图~\ref{fig:PhaseWidth}c，当 $\phi_\mathrm{a}=0.50$ 时，$f_\Omega^\ast$ 随 $\phi_{\mathrm A}^{\mathrm{iso}}$ 的变化趋势在不同 $N\chi$ 下表现出明显差异。
% 在较低 $N\chi$ 条件下，$f_\Omega^\ast$ 在 B 组分胶束中心趋近于 $0$，表明 L-S TT 更倾向于线型 diblock 构型；而在较高 $N\chi$ 条件下，胶束中心处的 $f_\Omega^\ast$ 则保持为非零的有限值，说明该区域出现了对星型 A$_2$AB 构型的偏好。
% 此外，在 A 组分作为 shell 的组成区域中，$f_\Omega^\ast$ 的取值也随 $N\chi$ 的增大而逐渐降低。
% 图~\ref{fig:PhaseWidth}d 进一步表明，在固定 $\phi_{\mathrm A}^{\mathrm{iso}}=0.75$ 时，最高 $N\chi$ 对应的 $f_\Omega^\ast$ 曲线近似为一条非零水平线，说明在强分离条件下，该区域的最优拓扑选择已基本摆脱对 $\phi_{\mathrm{a}}$ 的依赖，并锁定在一个稳定的非零值附近。
% 这些结果表明，在低 $N\chi$ 区域，体系的拓扑选择主要受滑动熵调控。
% 随着 $N\chi$ 增加，界面相关的焓效应逐步增强，并不断削弱这种由熵主导的拓扑择优，使 $f$ 的分布由以整体构象熵调节为主，逐渐转向以缓解局部 packing frustration 为主，进而表现为对 $\phi_\mathrm{a}$ 依赖的减弱，以及在不同空间位置（$\phi_{\mathrm A}^{\mathrm{iso}}$）上的功能分化增强。
% 需要指出的是，在本研究考察的 $N\chi\le 45.0$ 范围内，这种焓主导趋势尚未完全压制滑动自由度所带来的构象调节收益，因此 BCC 相仍能够持续从拓扑调节中获益，并保持稳定扩张的相区宽度。
To clarify how different topology-regulation mechanisms contribute to the formation of complex ordered structures, we examine the evolution of the BCC and FCC phase widths, $\Delta\bar{\phi}_{\mathrm A}$, with $N\chi$ and $\phi_{\mathrm a}$, as shown in Fig.~\ref{fig:PhaseWidth}b,e.
For all $\phi_{\mathrm a}$, the BCC phase width increases monotonically with $N\chi$, indicating that BCC consistently benefits from the conformational regulation enabled by the sliding degree of freedom.
This behavior can be understood from the distribution of the preferred topology $f_\Omega^\ast$ on different density layers $\phi_{\mathrm A}^{\mathrm{iso}}$.
As shown in Fig.~\ref{fig:PhaseWidth}c, for $\phi_{\mathrm a}=0.50$, the variation of $f_\Omega^\ast$ with $\phi_{\mathrm A}^{\mathrm{iso}}$ depends strongly on $N\chi$.
At lower $N\chi$, $f_\Omega^\ast$ approaches zero near the B-rich micellar center, indicating a preference for an AB linear diblock configuration.
At higher $N\chi$, however, $f_\Omega^\ast$ remains finite in the same region, showing that an A$_2$AB-like topology becomes locally favored.
In the A-shell-forming composition range, $f_\Omega^\ast$ also decreases gradually as $N\chi$ increases.
Fig.~\ref{fig:PhaseWidth}d further shows that, at fixed $\phi_{\mathrm A}^{\mathrm{iso}}=0.75$, the $f_\Omega^\ast$ curve at the highest $N\chi$ is nearly horizontal and remains at a finite value.
This indicates that, under strong segregation, the preferred topology in this region becomes largely independent of $\phi_{\mathrm a}$ and is locked around a stable nonzero value.
These results suggest a crossover in the role of sliding topology.

At low $N\chi$, the topology preference remains close to the RPA--MUT reference, reflecting the stronger role of the conformational response associated with the sliding ensemble.
As $N\chi$ increases, stronger interfacial and packing constraints weaken this global topology preference and promote increasingly local topology regulation in response to heterogeneous filling environments, as reflected by the reduced dependence on $\phi_{\mathrm a}$ and the stronger functional differentiation across spatial layers.

Within the range considered here, $N\chi\le45.0$, the stronger interfacial constraints do not eliminate the benefit of local topology adjustment.
Consequently, the BCC phase continues to gain stability from topology regulation and its phase width expands steadily.


% 相比 BCC 相，图~\ref{fig:PhaseWidth}e 显示 FCC 相区宽度 $\Delta \bar{\phi}_{\mathrm{A}}$ 随 $N\chi$ 呈现明显的非单调变化：在中等分离强度 $N\chi \approx 35$ 附近达到最大，而在更高 $N\chi$ 下则略有收缩。
% 这表明，滑动自由度对 FCC 相的稳定化作用只在有限的分离强度范围内最为有效；当界面相关焓贡献进一步增强时，这种拓扑调节收益反而开始减弱。
% 结合图~\ref{fig:PhaseWidth}f,g 中 $f_\Omega^\ast$ 的分布特征可以理解这一行为。
% 如图~\ref{fig:PhaseWidth}f 所示，在 $\phi_\mathrm{a}=0.50$ 时，FCC 相中 $f_\Omega^\ast$ 对空间位置 $\phi_{\mathrm A}^{\mathrm{iso}}$ 的响应主要集中在 A shell 对应的高 $\phi_{\mathrm A}^{\mathrm{iso}}$ 区域，而在其余区域内始终为 0，即维持 AB 线型 diblock 构型。
% 这说明，相较于 BCC 相中较为连续的空间调节，FCC 相中的滑动拓扑调节已被压缩到更狭窄的局域区域，其可参与构象优化的空间范围明显受限。
% 图~\ref{fig:PhaseWidth}g 表明，在固定 $\phi_{\mathrm A}^{\mathrm{iso}}=0.75$ 时，$f_\Omega^\ast$ 随 $\phi_\mathrm{a}$ 的变化同样呈现出与 BCC 相相似的 $N\chi$ 依赖行为。
% 因此，FCC 相区所表现出的极值行为可理解为熵效应与焓效应短暂平衡的结果。
% 在中等 $N\chi$ 下，滑动自由度仍能在界面附近提供有限但有效的拓扑重排，从而帮助体系适应 FCC 所要求的更复杂局域填充与曲率约束，因此 FCC 相区得以扩张。
% 而当 $N\chi$ 进一步增大时，界面约束和局域 packing frustration 显著增强，导致体系已难以通过进一步的拓扑分化来持续降低界面自由能，最终表现为 FCC 相稳定性下降及其相区宽度的回缩。
% 总体而言，BCC 与 FCC 的差异表明，滑动自由度对有序相的稳定化作用并非简单随界面复杂度单调增强，而是取决于该结构是否允许拓扑调节在足够大的空间范围内持续分担局域填充与界面匹配任务。
% 这一现象完整呈现了 L-S TT 体系中熵和焓两种机制之间的竞争和协同关系。
Compared with BCC, the FCC phase shows a nonmonotonic dependence of the composition width $\Delta\bar{\phi}_{\mathrm A}$ on $N\chi$ (Fig.~\ref{fig:PhaseWidth}e).
The FCC composition window reaches its maximum width near $N\chi\approx35$ and then slightly narrows at higher $N\chi$, while FCC remains stable at selected compositions up to $N\chi=45$ within the parameter range examined here. 


The contrasting BCC and FCC phase widths are accompanied by different patterns of topology preference across the density layers.
For soft polymer domains, the more nearly spherical Voronoi cell of BCC can reduce the geometric distortion required for nominally spherical domains to fill space.
Sliding provides an additional means of adapting chain conformations to the local filling requirements of each structure.


Fig.~\ref{fig:PhaseWidth}f,g resolves this topology response in FCC.
For $\phi_{\mathrm a}=0.50$, the response of $f_\Omega^\ast$ to the spatial layer $\phi_{\mathrm A}^{\mathrm{iso}}$ in FCC is mainly confined to the high-$\phi_{\mathrm A}^{\mathrm{iso}}$ region associated with the A shell, whereas $f_\Omega^\ast$ remains close to zero in the other regions, corresponding to the AB linear diblock configuration.
Compared with the more continuous variation of the preferred topology in BCC, this response in FCC is concentrated within a narrower range of density layers.
Fig.~\ref{fig:PhaseWidth}g further shows that, at fixed $\phi_{\mathrm A}^{\mathrm{iso}}=0.75$, the dependence of $f_\Omega^\ast$ on $\phi_{\mathrm a}$ follows an $N\chi$-dependent trend similar to that observed in BCC.
Together, these trends suggest that the maximum in the FCC phase width reflects a competition between local conformational adaptation and stronger interfacial constraints as $N\chi$ increases.
The entropy loss in Fig.~\ref{fig:SpherePOmega}c provides a quantitative comparison of the accompanying topology homogenization.
At $N\chi=45$ and $\phi_{\mathrm a}=0.50$, the representative A-shell states exhibit a larger sliding-distribution entropy loss in BCC than in FCC.
The smaller value for FCC indicates a weaker departure of the local topology distributions from uniform sliding under the A segment weighting, and hence stronger topology homogenization in this representative state.
The same BCC--FCC ordering is also observed for the representative A core states.
Redistribution of topology weights may help accommodate the local filling requirements of FCC at intermediate $N\chi$, whereas its subsequent narrowing indicates that competing phases become favorable over a larger composition range at higher $N\chi$.
The contrast with the steadily expanding BCC window suggests that the stabilizing effect of sliding depends on how effectively the accessible topology response adapts to the distinct local filling geometry of each structure.

